% =====================================================================
\chapter{Resolución de triángulos y producto escalar}
\label{tema:3}
% =====================================================================

Este tema extiende la trigonometría a triángulos cualesquiera mediante los teoremas del seno y del coseno, clasifica los casos de resolución de triángulos (incluido el caso ambiguo) y muestra que el teorema del coseno es la misma afirmación que la fórmula geométrica del producto escalar. Se sigue~\cite[§10.1, §10.2 y §10.8]{osat}, el tratamiento de los triángulos oblicuos de~\cite{stewart} y, para el producto escalar, la sección~1.2 de Strang~\cite{strang}.

\section{Área y teorema del seno}

En lo que sigue, $ABC$ es un triángulo con ángulos $A,B,C$ (que se identifican con sus medidas en radianes) y lados $a=BC$, $b=CA$, $c=AB$, de modo que cada lado es el opuesto al ángulo de la misma letra. Se cumple $A+B+C=\pi$ y todos los ángulos están en $(0,\pi)$, de modo que sus senos son positivos.

\begin{proposicion}[Área]\label{prop:area}
El área de $ABC$ es $S=\tfrac12 ab\sin C=\tfrac12 bc\sin A=\tfrac12 ca\sin B$.
\end{proposicion}

\begin{teorema}[Teorema del seno]\label{teo:seno}
En todo triángulo,
\[
\frac{a}{\sin A}=\frac{b}{\sin B}=\frac{c}{\sin C}.
\]
\end{teorema}

El valor común de los tres cocientes es el diámetro $2R$ de la circunferencia circunscrita (\cref{eje:3-circunradio}).

\section{Teorema del coseno}

\begin{teorema}[Teorema del coseno]\label{teo:coseno}
En todo triángulo, $c^{2}=a^{2}+b^{2}-2ab\cos C$, y análogamente para los otros dos lados.
\end{teorema}

\begin{corolario}\label{cor:pitagoras}
$C=\pi/2$ si y solo si $c^{2}=a^{2}+b^{2}$ (teorema de Pitágoras y su recíproco); $C$ es obtuso si y solo si $c^{2}>a^{2}+b^{2}$, y agudo si y solo si $c^{2}<a^{2}+b^{2}$.
\end{corolario}

\section{Resolución de triángulos}

\emph{Resolver} un triángulo es calcular sus seis elementos a partir de tres datos que incluyan al menos un lado. La \cref{tab:casos} resume qué teorema se usa en cada caso~\cite[§10.1--10.2]{osat}.

\begin{table}[ht]
\centering\small
\caption{Casos de resolución de triángulos.}\label{tab:casos}
\begin{tabular}{@{}>{\raggedright\arraybackslash}p{4.6cm}>{\raggedright\arraybackslash}p{4.6cm}>{\raggedright\arraybackslash}p{4.6cm}@{}}
\toprule
Datos & Primer paso & Número de soluciones \\
\midrule
LLL (tres lados) & teorema del coseno para un ángulo & una, si se cumple la desigualdad triangular \\
LAL (dos lados y el ángulo comprendido) & teorema del coseno para el tercer lado & una \\
ALA o AAL (dos ángulos y un lado) & $C=\pi-A-B$ y teorema del seno & una, si $A+B<\pi$ \\
LLA (dos lados y un ángulo no comprendido) & teorema del seno & cero, una o dos (caso ambiguo) \\
\bottomrule
\end{tabular}
\end{table}

\begin{proposicion}[Caso ambiguo]\label{prop:ambiguo}
Se dan $a$, $b$ y el ángulo $A$ opuesto a $a$, con $A$ agudo, y sea $h=b\sin A$. Entonces existen: ningún triángulo si $a<h$; uno (rectángulo en $B$) si $a=h$; dos si $h<a<b$; y uno si $a\ge b$.
\end{proposicion}

\addtocounter{definicion}{1}%

Cuando se conocen los tres lados, el área puede calcularse sin ángulos mediante la \emph{fórmula de Herón}, $S=\sqrt{s(s-a)(s-b)(s-c)}$ con $s=\frac{a+b+c}{2}$, que se demuestra en el \cref{eje:3-heron}.

\section{Producto escalar y ángulo entre vectores}

\begin{definicion}[Producto escalar]
Para $u=(u_1,\dots,u_n)$ y $v=(v_1,\dots,v_n)$ en $\R^{n}$, el \emph{producto escalar} es $u\cdot v=\sum_{i=1}^{n}u_iv_i$ y la \emph{norma} es $\norm{u}=\sqrt{u\cdot u}$.
\end{definicion}

El producto escalar es simétrico y bilineal; en particular,
\begin{equation}\label{eq:norma-diferencia}
\norm{u-v}^{2}=(u-v)\cdot(u-v)=\norm{u}^{2}+\norm{v}^{2}-2\,u\cdot v .
\end{equation}

\begin{teorema}[Forma geométrica del producto escalar]\label{teo:escalar}
Sean $u,v\in\R^{2}$ (o $\R^{3}$) no nulos y $\varphi\in[0,\pi]$ el ángulo que forman. Entonces
\[
u\cdot v=\norm{u}\,\norm{v}\cos\varphi .
\]
\end{teorema}

\begin{corolario}\label{cor:escalar}
Para $u,v$ no nulos del plano o del espacio: (i) $u\perp v$ si y solo si $u\cdot v=0$; (ii) $\abs{u\cdot v}\le\norm{u}\norm{v}$ (desigualdad de Cauchy--Schwarz); (iii) la componente de $u$ en la dirección de $v$ es $\norm{u}\cos\varphi=u\cdot v/\norm{v}$, y la proyección de $u$ sobre la recta de $v$ es $\frac{u\cdot v}{\norm{v}^{2}}\,v$.
\end{corolario}

En $\R^{n}$ con $n>3$ no hay triángulos que dibujar y se procede al revés: la desigualdad de Cauchy--Schwarz se demuestra algebraicamente y se \emph{define} el ángulo por $\cos\varphi=u\cdot v/(\norm{u}\norm{v})$ (módulo~2, tema~2). El cociente recibe el nombre de \emph{similitud coseno}.

\addtocounter{definicion}{1}%

\addtocounter{definicion}{1}%

\addtocounter{definicion}{1}%

